\subsection{Size of the two-mark Pfaffian denominator}\label{c4:pair-denominator}

The quotient in \eqref{c4:Pair-sum} contains the square of the vacuum
normalization. Its contribution is a power of the total number of sites,
which will cancel the excess radial growth when the marked arcs have
different lengths. We obtain that power from a positive integral
representation of a ratio of Pfaffians.

\paragraph{The exact formula and a comparison Pfaffian.}
For even \(N\), the empty component of the \(p=1\) vacuum is
\[
 D_1(X)=\prod_{i<j}\frac{p_K(x_i,x_j)}{x_j-x_i}
 \operatorname{Pf}\left[\frac{x_j^2-x_i^2}{p_K(x_i,x_j)}\right]_{i,j=1}^N,
 \qquad p_K(x,y)=x^2+\sqrt2xy+y^2.
\]
Here \(x_i=t_i\) in the vacuum notation. Skew alternation cancels the
Vandermonde denominator, and each possible \(p_K\)-pole is canceled by
the outside product, so this is a symmetric polynomial of the required
degree. At \(x_j=q^3x_i=q^3u\), only the Pfaffian term contracting these
two sites survives. With the pair adjacent, its prefactor is
\((1+q^3)u\prod_{k\ne i,j}(x_k-q^{-3}u)(x_k-q^6u)\), exactly the
empty-component recursion. Symmetry gives the reversed pair ratio as
well. These values determine the polynomial by interpolation in any
site, starting with \(D_1(x_1,x_2)=x_1+x_2\). Put
\[
 \widetilde D(X)=D_1(X)\prod_i x_i^{-(N-1)/2}.
\]
Inverting all arguments preserves \(\widetilde D\), with consistent
logarithms; the signs from the outside product and the Pfaffian cancel
because \(N\) is even.

Take \(x_j=e^{i\theta_j}\), with all angles in a small fixed neighborhood
of an interval of width \(\pi/4\). Fix that neighborhood throughout the
argument. Its width \(\Theta\) is less than \(\pi b'\), where
\(b'=3/4\), so after centering the angles the numbers
\(p_j=e^{i\theta_j/b'}\) satisfy \(\Re p_j\ge a>0\) with a fixed
\(a\). The estimates below are uniform over these arrays, without a
minimum separation condition. Define the continuous function
\[
 g(s)=\frac{(2\cos s+\sqrt2)\tan(s/(2b'))}{2\sin(s/2)},
 \qquad g(0)=\frac{2+\sqrt2}{2b'}.
\]
We will prove
\begin{equation}
 |\widetilde D(X)|^2\ge
 (2b')^{-N}\prod_{i<j}g(\theta_i-\theta_j)^2\,
 N^{5/24-o(1)}.                                      \label{c4:Pair-Pf}
\end{equation}
The error depends on the fixed angular neighborhood, but not on the
array or its separation. We first use distinct angles and pass to
collisions at the end.

For \(w=\log(x_j/x_i)\), the Pfaffian kernel above has the representation
\[
 \frac{\sinh w}{\cosh w+1/\sqrt2}
 =2\int_0^\infty \sin(wT)
             \frac{\cosh(\pi T/4)}{\sinh(\pi T)}\,dT.
\]
One can verify it for real \(w>0\) by integrating the full exponential
in principal value in the upper half-plane, with upper edges at
half-integer heights. The residues give
\(i(1+2\sum_{n\ge1}(-1)^n\cos(n\pi/4)e^{-nw})\); analytic
continuation gives the stated formula on the required strip.
Replacing the density \(\cosh(\pi T/4)/\sinh(\pi T)\) by
\(1/(2\sinh(\pi b'T))\) gives the comparison kernel
\[
 K_0(w)=\frac1{2b'}\tanh\!\left(\frac{w}{2b'}\right).
\]
Its Pfaffian is
\[
 \operatorname{Pf}[K_0(\log(x_j/x_i))]
 =(2b')^{-N/2}\prod_{i<j}
          \tanh\!\left(\frac{\log(x_j/x_i)}{2b'}\right).
\]
This is the Schur skew-Cauchy identity: induction compares the residues
at \(x_i^{1/b'}=-x_j^{1/b'}\), and alternation removes the remaining
constant. Substitution into the outside product defining
\(\widetilde D\) gives a comparison magnitude whose square is exactly
\((2b')^{-N}\prod_{i<j}g(\theta_i-\theta_j)^2\).

\paragraph{A positive ratio and its counting density.}
Order the distinct angles once, and let \(\mathcal R_N\) be the ratio
of the original kernel Pfaffian to the comparison kernel Pfaffian.
Expanding the comparison Pfaffian by wedging its two-vector integral
kernels gives a probability law of \(N/2\) variables \(T_k>0\), with
density proportional to
\[
 \det[e^{\pm\theta_jT_k}]_{j;(k,\pm)}
          \prod_k\frac{dT_k}{\sinh(\pi b'T_k)}.
\]
Choose the constant sign of the determinant to make this density
positive. Its sign is constant: after ordering the exponent columns,
total positivity follows because a nonzero linear combination of
\(r\) distinct real exponentials has at most \(r-1\) distinct zeros
(divide by one exponential, differentiate, and induct). Permuting
entire column pairs has even sign. Alternation within each pair gives
integrability at \(T_k=0\); the angular width \(\Theta<\pi b'\)
gives integrability at infinity. The same fixed phase occurs in both
Pfaffians and cancels in their ratio. Thus
\[
 \mathcal R_N=\mathbb E\prod_{k=1}^{N/2}r(T_k)>0,
 \qquad
 r(T)=\frac{(1+e^{-u})(1-e^{-3u})}{1-e^{-4u}}
     =1+\frac1{2\cosh u}\ge1,
 \quad u=\pi T/2.
\]
Define the one-point counting density \(\rho_N\) by
\(\mathbb E\sum_k f(T_k)=\int_0^\infty f(T)\rho_N(T)\,dT\).
Jensen's inequality now reads
\begin{equation}
 \log\mathcal R_N\ge
        \int_0^\infty\rho_N(T)\log r(T)\,dT.             \label{c4:Pair-Jensen}
\end{equation}
We will show that for each fixed \(T>0\),
\[
 \rho_N(T)=\frac{b'}\pi\log N+O_T(1),
\]
uniformly over distinct arrays. This concerns the density at bounded
\(T\); the law itself always has exactly \(N/2\) variables.

Write
\[
 K_{ij}=\frac{p_i-p_j}{p_i+p_j},\qquad
 w_j(y)=\frac{e^y-p_j}{e^y+p_j},\qquad B(y)=\prod_jw_j(y),
 \qquad k=b'T.
\]
The skew integrals of
\(f_i^+f_j^--f_i^-f_j^+\), where
\(f_j^\pm=e^{\pm\theta_jT}=p_j^{\mp ik}\), against
\(dT/\sinh(\pi b'T)\) form the matrix \(K/(ib')\).
Taking the half log-determinant derivative of this Pfaffian therefore
gives the exact density formula
\[
 \rho_N(T)=\frac{ib'}{\sinh(\pi k)}
              (p^{ik})^tK^{-1}p^{-ik}.
\]
The matrix is invertible for distinct angles in the fixed interval.
Its inverse is controlled through the skew Schur complement in the
same Pfaffian identity:
\[
 w(y)^tK^{-1}w(y')=
       h(y-y')[B(y)B(y')-1],\qquad h(s)=\tanh(s/2).
\]
Also, contour shift followed by residues or the Euler integral gives
\[
 \int_{\mathbb R}e^{iky}w_j'(y)\,dy=c(k)p_j^{ik},
 \qquad c(k)=\frac{2\pi k}{\sinh(\pi k)}.
\]
These two identities reduce the density estimate to a uniform bound
for a mixed derivative.

\paragraph{The logarithmic local density.}
Put \(L_N=\log N\), \(I_N(y)={\bf1}_{\{|y|<L_N\}}\), and
\(J_N=1-I_N\). Uniformly in the site array,
\[
 |B(y)|\le e^{-cNe^{-|y|}},\qquad
 |B(y)-1|\le CNe^{-|y|}\quad (|y|>L_N),\qquad
 \|B'\|_1=O(1).
\]
The second estimate uses even \(N\) at the negative end. For the third,
product differentiation gives
\(|B'(y)|\le CNe^{-|y|}e^{-c(N-1)e^{-|y|}}\).
Integrating the first bound on the central interval and the second on
its complement yields
\(\|B-J_N\|_1=O(1)\). All constants use only the fixed lower bound
on \(\Re p_j\).

For \(F(y,y')=h(y-y')[B(y)B(y')-1]\), differentiation gives
\[
 \begin{aligned}
 \partial_y\partial_{y'}F
 ={}&h''(y-y')[1-B(y)B(y')]\\
 &+h'(y-y')[B(y)B'(y')-B'(y)B(y')]\\
 &+h(y-y')B'(y)B'(y').
 \end{aligned}
\]
The last two terms have bounded \(L^1(dy\,dy')\) norm, by
\(|B|\le1\), \(\|B'\|_1=O(1)\), and integrability of \(h'\).
In the first term replace both factors \(B\) by \(J_N\), at bounded
\(L^1\) cost. The difference between the result and
\(h''(y-y')I_N(y)\) is
\(h''(y-y')J_N(y)I_N(y')\); it too has bounded \(L^1\) norm because
\(h''\) has a finite first absolute moment. Consequently
\[
 \partial_y\partial_{y'}F
     =h''(y-y')I_N(y)+\mathcal E_N(y,y'),\qquad \|\mathcal E_N\|_1=O(1).
\]
Integrate against \(e^{ik(y-y')}\). The interval selected by \(I_N\)
has length \(2\log N\), and
\(\int e^{iks}h''(s)\,ds=-ikc(k)\). The derivative of the left-hand
Schur-complement expression therefore gives
\[
 c(k)^2(p^{ik})^tK^{-1}p^{-ik}
       =-2ikc(k)\log N+O(1).
\]
Substitution in the density formula proves the asserted
\(\rho_N(T)/\log N\to b'/\pi\), uniformly for each fixed \(T>0\).

Finally \(\rho_N\log r\ge0\), so Fatou applied to
\eqref{c4:Pair-Jensen} gives, along any sequence of distinct site arrays,
\[
 \liminf_{N\to\infty}\frac{\log\mathcal R_N}{\log N}
 \ge\frac{b'}\pi\int_0^\infty\log r(T)\,dT=\frac5{48}.
\]
For the integral expand the three logarithms in the displayed formula
for \(r\): its value is
\(\frac2\pi(\pi^2/12-\pi^2/18+\pi^2/24)=5\pi/36\).
Since the argument applies to every sequence of arrays, the lower
bound has a uniform \(o(1)\). Squaring the ratio and restoring the
comparison magnitude proves \eqref{c4:Pair-Pf} for distinct sites.
For each fixed \(N\), continuity of \(\widetilde D\) and of the product
of \(g\)'s then proves the same bound at colliding sites. No inverse
matrix needs to be continued through a collision.
